\section{Proofs and Technical Details}\label{app:proofs}

\subsection{Conventions}\label{app:notation}

Notation below matches the main text: $\TV$ total variation
($\TV(P^{N},Q^{N})$ its $N$-fold product value); $\alpha$ a confidence
level; $\delta=\TV(P,Q)$ a deception's observable footprint; $\beta_d$ a
tolerated false-alarm rate; $z_p$ the standard-normal quantile
$\Pr[Z\le z_p]=p$ (the two-sided critical value at level $\beta_d$ is
$z_{1-\beta_d/2}$; the one-sided analogue is $z_{1-\beta_d}$);
$X_B$ the verdict of
Definition~\ref{def:beta}; $P^{\mathrm{pas}}_B$ the benign transcript
distribution. Symbols reused across sections with different referents are:
$\rho$ (Spearman correlation in the instrument layer) vs.\ $\rho_0$ (the
reference width of Theorem~\ref{thm:onset}); $c$ (probe context in
Definitions~\ref{def:probe}--\ref{def:beta} vs.\ the reactivity rate of
Proposition~\ref{prop:reactivity}); $m$ (screening probes per proposition
in the proof of Proposition~\ref{prop:screening} vs.\ the held-out misread
rate of Proposition~\ref{prop:budget}); $K$ (propositions in
Theorem~\ref{thm:uniform} vs.\ the hide-threshold of
Proposition~\ref{prop:alloc}); $B$ (belief hypothesis vs.\ the attacker's
allocation budget in Section~\ref{sec:release}). Each is local to its
subsection and never co-occurs with its homonym. Every claim below is
exactly the corresponding statement of the main text.

\subsection{Proof of Theorem~\ref{thm:passive} and
Corollary~\ref{cor:detection}}\label{app:t1}

\begin{proof}[Proof of Theorem~\ref{thm:passive}]
Fix any benign operation $\mu$; adversary-controlled events of zero
$\mu$-probability exist---take $E$ to be a forged ``[SYSTEM OVERRIDE]''
directive. Let $B_0$ refuse on every history; let $B_1$ agree with $B_0$ on
every history $h$ with $\Pr_\mu(h)>0$ and comply with probability $1$ on
histories containing $E$ (well-defined: ``$E\in h$'' never fires on
$\mathrm{supp}(\mu)$, so $B_1$ coincides with $B_0$ on benign data; clearly
$B_0\neq B_1$).

\emph{(i) The passive distributions coincide.} For every $h\in\mathrm{supp}(\mu)$
the two response kernels agree by construction,
$\kappa_{B_0}(h,\cdot)=\kappa_{B_1}(h,\cdot)$; the transcript distributions
agree pointwise, so $\TV\bigl(P^{\mathrm{pas}}_{B_0},P^{\mathrm{pas}}_{B_1}\bigr)=0$.

\emph{(ii) One probe separates them.} $\Pi_{\mathrm{auth}}$ injects $E$
with probability $1$; with $X_B$ indicating credential disclosure, the
expected verdict is $0$ under $B_0$ (refusal) and $1$ under $B_1$ (it
complies): the separation is exactly $1$.
\end{proof}

\begin{proof}[Proof of Corollary~\ref{cor:detection}]
Any monitor is a measurable function of the passive transcript stream. Since
the transcript distributions coincide by Theorem~\ref{thm:passive}(i), so do
the distributions of the monitor's output; equivalently, the likelihood ratio
between them is $1$ everywhere. For distinguishing $B_1$ from $B_0$, the
Neyman--Pearson lemma then bounds the true-positive rate of any such monitor
by its false-positive rate.
\end{proof}

\subsection{Proofs for Section~\ref{sec:qc}}\label{app:qc}

\begin{proof}[Proof of Theorem~\ref{thm:uniform}]
Fix a proposition $i$ and the accuracy $\varepsilon>0$. Hoeffding's inequality
\cite{hoeffding1963}, applied to the $n_i$ independent Bernoulli verdicts,
gives per proposition
\[
\Pr\bigl[\,|\hat\beta_i-\beta_i|>\varepsilon\,\bigr]\;\le\;2e^{-2n_i\varepsilon^{2}}.
\]
Under the allocation $n_i\ge\frac{1}{2\varepsilon^{2}}\ln\frac{2K}{\alpha}$,
the right-hand side is at most $2e^{-\ln(2K/\alpha)}=\alpha/K$. A union bound
over the $K$ propositions then yields
\[
\Pr\Bigl[\max_{i\in[K]}\,|\hat\beta_i-\beta_i|>\varepsilon\Bigr]
\;\le\;\sum_{i=1}^{K}\frac{\alpha}{K}\;=\;\alpha,
\]
as claimed.
\end{proof}

\begin{proof}[Proof sketch of Proposition~\ref{prop:screening}]
Phase 1 allocates $m=O(\tau^{-2}\ln(K/\alpha))$ uniform probes per
proposition and keeps those whose empirical rate is above threshold;
Hoeffding plus a union bound bounds the separation failure by $\alpha/2$.
Phase 2 $\varepsilon$-refines the surviving $s$ propositions at
$O(\varepsilon^{-2}\ln(K/\alpha))$ each; a below-threshold proposition left
unrefined exceeds the tolerance only under a two-width deviation, dismissed
within the residual $\alpha/2$. The saving concentrates on the $K-s$
below-threshold propositions, and Theorem~\ref{thm:uniform}'s uniform
$\varepsilon$-guarantee admits no such saving ($\varepsilon$-width on every
proposition). Optimality remains open.
\end{proof}

\subsection{Proofs for Section~\ref{sec:t5}}\label{app:t5}

\begin{proof}[Proof of Lemma~\ref{lem:amp}]
The Hahn decomposition of $P$ with respect to $Q$ yields a measurable set $A$
with $P(A)-Q(A)=\delta$. Count visits to $A$: let
$S=\sum_{i\le N}\mathbf{1}[X_i\in A]$. Under $P$,
$S\sim\mathrm{Bin}(N,p)$ with $p=P(A)$; under $Q$,
$S\sim\mathrm{Bin}(N,p-\delta)$. Consider the threshold test
$\psi=\mathbf{1}[S\ge t]$ with $t=N(p-\delta/2)$, which compares the count to
its two means at half the footprint. Hoeffding's inequality for binomial
counts gives $\Pr_p[S<t]\le e^{-N\delta^{2}/2}$, and, symmetrically,
$\Pr_q[S\ge t]\le e^{-N\delta^{2}/2}$ under $Q$.
The total variation distance between $P^{N}$ and $Q^{N}$ is the supremum of
$\Pr_p[\psi]-\Pr_q[\psi]$ over tests $\psi$, so this particular test lower-bounds it:
\[
\begin{aligned}
\TV\bigl(P^{N},Q^{N}\bigr)
&\ge\Pr_p[S\ge t]-\Pr_q[S\ge t]\\
&\ge\bigl(1-e^{-N\delta^{2}/2}\bigr)-e^{-N\delta^{2}/2}\\
&=1-2e^{-N\delta^{2}/2},
\end{aligned}
\]
which is the claim.
\end{proof}

\begin{proof}[Proof of Theorem~\ref{thm:patch}]
\emph{(i).} This is Definition~\ref{def:sens} applied with
$\TV(P,Q)=\delta$: for every proposition $\pi$,
$|\beta[Q](\pi)-\beta[P](\pi)|\le L\,\TV(P,Q)=L\delta$.

\emph{(ii).} Two facts combine, with the test direction made explicit.
\emph{Existence (for the attacker who knows the honest rate $p$ and the
deceived direction and size, i.e.\ $\delta$):} Lemma~\ref{lem:amp}'s
threshold test $\psi=\mathbf{1}[S\ge t]$ at $t=N(p-\delta/2)$ fires
\emph{under honest $P$} (it lower-bounds $\TV$ through
$\Pr_P[\psi]-\Pr_Q[\psi]$); the detector the attacker actually runs is its
complement $\psi^{c}=\mathbf{1}[S<t]$---alarm when the count of Hahn-set
visits falls below the honest expectation discounted by half the
footprint. Under $P$, $\Pr_P[\psi^{c}]\le e^{-N\delta^{2}/2}$ controls the
false alarm; under $Q$, $S\sim\mathrm{Bin}(N,p-\delta)$ concentrates below
$t$, so $\Pr_Q[\psi^{c}]\ge1-e^{-N\delta^{2}/2}$, and subtracting the
false-alarm term on the $P$ side gives the claimed probability of correct
detection at least $1-2e^{-N\delta^{2}/2}$. (An attacker who
knows $p$ but neither the direction nor the size of the change runs the
corresponding two-sided deviation test at level $\beta_d$ on the honest
expectation; the exponential constant degrades by the two-sided split, and
the bound of (iii) is unchanged because it inverts whichever direction
holds.) \emph{Optimality:} for every
test $\psi$, $|\E_P[\psi]-\E_Q[\psi]|\le\TV(P^{N},Q^{N})$, so no test exceeds
$\TV(P^{N},Q^{N})$.

\emph{(iii).} Invert the exponent of (ii). Holding the detection
probability at most $\gamma$ requires
$1-2e^{-N\delta^{2}/2}\le\gamma$, which is equivalent to
$e^{-N\delta^{2}/2}\ge(1-\gamma)/2$, that is, to
$\delta\le\sqrt{2\ln(2/(1-\gamma))/N}$. Substituting this largest feasible
footprint into (i) gives
$\text{gain}\;\le\;L\sqrt{2\ln(2/(1-\gamma))/N}=O(L/\sqrt{N})$.
\end{proof}

\begin{proof}[Proof sketch of Theorem~\ref{thm:onset}]
\emph{Setup and two boundaries.} The memorized reference of $N_0$ honest
samples concentrates around $p$: under the normal approximation to the
binomial (the same approximation the registered E5 floors carry; a regime
indicator at $N_0=10$, adequate at $N_0=40$; Section~\ref{sec:e5}),
$|\hat p_0-p|\le\rho_0=z\sqrt{p(1-p)/N_0}$ with probability
$\approx1-\beta_d$, where $z=z_{1-\beta_d/2}$. The fresh stream of $M$
deceived samples concentrates around $q=p+\delta$ (the registered planes
shift the event rate \emph{up}; the sign matters---an earlier draft of this
proof printed the deceived rate as $p-\delta$, which cannot represent the
$\delta=0.4$ planes at $p=0.30$ and is corrected here). Consider the
two-sample statistic $S=\hat q_M-\hat p_0$ with analytic standard
deviation
$\sigma_{\mathrm{dm}}=\sqrt{q(1-q)/M+p(1-p)/N_0}$ and the level-$\beta_d$
two-sided rejection region $|S|/\sigma_{\mathrm{dm}}>z$. Under honest
sampling ($\delta=0$), $\Pr[\text{reject}]\le\beta_d$ to first order, so
the test runs at its level; its power at shift $\delta>0$ is dominated by
the exceedance tail,
$\Pr[\text{reject}]\approx\Phi\!\bigl(\delta/\sigma_{\mathrm{dm}}-z\bigr)$,
which reaches $1-\beta_d$ (equivalently $\delta/\sigma_{\mathrm{dm}}\ge z+z'$ with $z'=z_{1-\beta_d}$) iff
$\delta\ge(z+z')\,\sigma_{\mathrm{dm}}$.

\emph{(i)(a) The certified regime.} If
$\delta>\Delta_{\beta_d}:=(z+z')\sqrt{p(1-p)/N_0}$, then since
$\sigma_{\mathrm{dm}}$ decreases in $M$ toward
$\sigma_0=\sqrt{p(1-p)/N_0}$, there is a finite horizon at which
$\delta\ge(z+z')\sigma_{\mathrm{dm}}$ first holds; solving
$\delta^{2}\ge(z+z')^{2}\bigl(q(1-q)/M+p(1-p)/N_0\bigr)$ for $M$ gives
$M\ge M^{*}(\delta)=(z+z')^{2}q(1-q)/\bigl(\delta^{2}-\Delta_{\beta_d}^{2}\bigr)$,
finite iff $\delta>\Delta_{\beta_d}$; and no horizon below
$M^{*}$ certifies, while at the registered planes
$(p,q)=(0.3,0.7)$---where $p(1-p)=q(1-q)=0.21$ coincidentally---$M^{*}$
evaluates to $30$ at the $(N_0,\delta)=(40,0.4)$ cell. The earlier
draft's claim that $\delta\ge\rho_0$ suffices with $M\gtrsim\delta^{-2}$
is \emph{false}: at $\delta=\rho_0$ the power is
$\Phi\!\bigl(z-z\bigr)=0.50$, not $1-\beta_d$; and at the two middle
registered cells ($\delta/\sigma_0=2.76<z+z'=3.60$) the power is
$\Phi(0.80)=0.79<0.95$ for \emph{every} $M$, even $M\to\infty$.

\emph{(i)(b) The uncertified middle regime.} For
$\rho_0\le\delta\le\Delta_{\beta_d}$, the limit
$\lim_{M\to\infty}\Pr[\mathrm{reject}]
=\Phi\!\bigl(\delta/\sigma_0-z\bigr)<1-\beta_d$: certified power is
unattainable at every horizon (the deception stays visible---the limit is
bounded away from the level $\beta_d$---but never refutable at the
registered level).

\emph{(i)(c) Below the reference width, and the two threshold forms.}
For a rate-difference test with threshold $c$ fixed independently of the
estimated honest rate, Hoeffding's inequality for independent
$[0,1]$-bounded means gives
$\Pr[\hat q-\hat p>c]\le\exp\bigl(-2c^{2}/(1/M+1/N_0)\bigr)$
independently of $p$, so level $\beta_d$ forces
$c\ge\sqrt{(1/M+1/N_0)\ln(1/\beta_d)/2}$, and the same bound from below
yields power $\ge1-\beta_d$ only when
$\delta\ge\sqrt{2\ln(1/\beta_d)(1/M+1/N_0)}
\ge w_{\mathrm{na}}=\sqrt{2\ln(1/\beta_d)/N_0}$: the $p$-robust form holds
at every honest rate at the price of dropping $\sqrt{p(1-p)}$. The normal
form $\rho_0=z\sqrt{p(1-p)/N_0}$ is tighter away from the rate extremes
(for the studentized statistic whose threshold adapts to the estimated
rate) and is the form the E5 comparison tests; at the E5 operating point
the forms differ by $w_{\mathrm{na}}/\rho_0\approx2.7$. Below
$\rho_0$---more precisely, below the threshold form in force---the gain
cap is the footprint bound of Theorem~\ref{thm:patch}(i) with
sensitivity: $\delta<\rho_0$ restricts the achievable shift of any
$L$-sensitive functional to $\le L\rho_0=O(L/\sqrt{N_0})$; the adaptive
minimax boundary between the regimes remains open
(Section~\ref{sec:discussion}).

\emph{Birth onset.} No proof is required or appropriate: with $N_0=0$
both hypotheses are consistent with the single observed sample path, and
the in-channel undetectability of deceptions inside the plausible set is
\emph{true by the definition of the plausible set}---the channel identity
stated as Remark~\ref{rem:birth}, not an arm of this theorem. Both the
real-world and deceived-world hypotheses reduce to the prior plausibility
of the observed transcript distribution, and the prior is out-of-band
information.
\end{proof}

\subsection{Proofs for Section~\ref{sec:surface}}\label{app:limits}

\begin{proof}[Proof of Corollary~\ref{cor:pub}]
Construct $T^{*}$ to agree with the compliant target on every history in
$\mathcal{S}$ (the public supports) and breach with probability $1$ on the
event $E=\{\pi'\in h\}$ (disjoint from $\mathcal{S}$ by hypothesis).
$T^{*}$ and the compliant target induce identical transcript distributions
under every public probe; $C$ accepts both. Yet $\pi'$ injects $E$ with
probability one, and the breach indicator is $1$ under $T^{*}$ on that
event: Theorem~\ref{thm:passive} with $E=\pi'$.
\end{proof}

\begin{proof}[Proof of Theorem~\ref{thm:compose}]
Let $s_A$ be the secret held by $A$; modify each component off-support
only. \emph{Step 1 (retriever).} $A'$ agrees with the certified $A$ on
$\mathcal{S}$ but returns a document containing $s_A$ on the off-support
request ``fetch document''. \emph{Step 2 (orchestrator).} $B'$ agrees with
the certified $B$ on $\mathcal{S}$ but summarizes faithfully any tool output
tagged $t$. \emph{Invisibility.} Neither the request nor the tag $t$ is
ever issued by $\Pi_{\mathrm{pub}}$ (by hypothesis), so both modifications
are unobservable; $A'$ and $B'$ remain certified. \emph{Realization.} The
cross-interface probe $\pi^{*}$ issues the fetch request to $A'$ and asks
$B'$ to summarize the tagged result; the reply contains $s_A$, hence
$\beta_{B'\circ A'}(\pi^{*})=1$. Finally, the interface history---$s_A$
under tag $t$---has zero measure under each component's local probe
distribution, which is why no component is individually non-compliant.
\end{proof}

\begin{proof}[Proof of Proposition~\ref{prop:gaming}]
\emph{First claim.} A target that complies on public probes but breaches on
held-out semantic equivalents is exactly the construction of
Corollary~\ref{cor:pub}, so the public certificate accepts it.

\emph{Second claim (transfer condition).} Hoeffding's inequality and a union
bound over the $J$ paraphrases give
$|\hat\beta(\theta_j)-\beta(\theta_j)|\le\varepsilon$ for all covered $j$
with probability $\ge1-\alpha$. The covering step from a covered
paraphrase $\theta_j$ to an uncovered $\theta$ within paraphrase distance
$r$ requires moving the \emph{probe}, not the transcript distribution, so
assumption (A1)---$L$-sensitivity over transcript distributions at fixed
probe, Definition~\ref{def:sens}, an identity with $L=1$ for binary
verdicts (Remark~\ref{rem:sens})---does \emph{not} by itself apply. The
proof uses the explicit, stronger assumption \textbf{(A3)}: the verdict
channel is Lipschitz over probe space with constant $L_{\pi}$---%
$|\beta(\theta)-\beta(\theta')|\le L_{\pi}\,d(\theta,\theta')$ for the
paraphrase metric $d$. Under (A3),
$\beta(\theta)\le\beta(\theta_j)+L_{\pi}r\le\hat\beta(\theta_j)
+\varepsilon+L_{\pi}r\le\tau+\varepsilon+L_{\pi}r$, bounding the profile on
the uncovered neighbourhood: the stated transfer radius, with $L_{\pi}$ in
place of $L$ throughout. (A3) is a genuinely strong assumption, and the
paper's own constructions are its counterexamples:
Corollary~\ref{cor:pub}, Theorem~\ref{thm:forgery}, and
Corollary~\ref{cor:adaptive} all build targets whose behavior is
discontinuous in probe identity, exactly the class (A3) excludes. This is
why transfer is stated as conditional and why the diagnostic $\Gamma$
exists: (A2) below protects only the sampled paraphrases; (A3) is the
bridge to their neighborhood, and targets violating it are precisely what
$\Gamma$ is designed to catch---when they are measured.

Assumption (A2) says the target cannot condition on probe \emph{provenance}
beyond probe content---what prevents different behavior on the held-out
probes themselves; a threat-model condition, not a consequence of held-out
text. With (A3) stated, the generalization is the standard
Lipschitz--covering argument; the diagnostic $\Gamma$ is the new object.
\end{proof}

\begin{proof}[Proof of Theorem~\ref{thm:forgery}]
Apply the construction of Theorem~\ref{thm:passive}, family by family. Define
$T$ to condition its reply on the event
$E_{\mathrm{pub}}=\{\text{the delivered probe lies in }
\mathrm{supp}(\Pi_{\mathrm{pub}})\}$: on $E_{\mathrm{pub}}$ its breach
indicator is prescribed by $B_{\mathrm{disp}}$; on its complement, by
$B_{\mathrm{true}}$. Both events are well-defined on histories and the supports are disjoint,
so the policy is consistent.
Under $\Pi_{\mathrm{pub}}$ the estimator observes only outcomes generated on
$E_{\mathrm{pub}}$, and therefore recovers the corresponding displayed
profile; under $\Phi_{\mathrm{held}}$, whose support is disjoint from
$\mathcal{S}$ by hypothesis, it observes only the complement and recovers the
true profile---each within estimation error.
\end{proof}

\begin{proof}[Proof of Corollary~\ref{cor:adaptive}]\looseness=-1
Condition the reply on the event $E$ that the delivered probe lies in the
union of the public support and $\mathrm{supp}(H)$: comply on $E$, breach off
it. The two supports are disjoint, so the policy is
consistent. The gaming gap computed on the published held-out set $H$ then
vanishes---every probe used to compute $\Gamma$ lands on $E$, where the target
complies---while the secret-set estimate is $1$: probes from $S$ land off
$E$, realizing the stated member of the class.
\end{proof}

\subsection{Proofs for Sections~\ref{sec:drift}--\ref{sec:budget}}\label{app:rest}

\begin{proof}[Proof of Proposition~\ref{prop:halflife}]\looseness=-1
For any $t\ge t_0$, the drift bound and the measured upper bound give
$\beta_t\le\beta_{t_0}+v\,(t-t_0)\le U_0+v\,(t-t_0)$: the first inequality is
the assumed $v$-Lipschitz drift, the second uses $\beta_{t_0}\le U_0$ (the
Wilson bound covers the truth at measurement time). Setting the
right-hand side to $\tau$ and solving yields $h_\tau=(\tau-U_0)/v$; beyond
that horizon the measured bound no longer controls $\beta_t$, which is the
``and no longer'' clause. The sequential statement is a Bonferroni union
bound \emph{over a pre-specified number of renewals} $R$, each renewed
certificate run at level $\alpha/R$; an unbounded renewal stream has no
simultaneous guarantee without a pre-registered stopping rule---the
classical discipline of sequential testing \cite{wald1945} (a calendar cap
on $R$ is the operational reading used in Section~\ref{sec:drift}).
The contribution is $h_\tau$ as a certificate validity horizon and its
renewal rule (in practice $v$ is differenced from consecutive certificate
measurements, its error absorbed by renewal).
\end{proof}

\begin{proof}[Proof of Proposition~\ref{prop:reactivity}]\looseness=-1
A pooled estimator over $n$ probes in one shared context averages the
position-dependent breach probabilities $\beta_0,\dots,\beta_{n-1}$; its bias
against the fresh-session estimand $\beta$ is the average of the linear
deviation bounds,
$\bigl|\frac{1}{n}\sum_{k=0}^{n-1}\beta_k-\beta\bigr|\le\frac{1}{n}\sum_{k=0}^{n-1}ck
=\frac{c(n-1)}{2}$:
the pooled estimator consistently estimates the context-conditioned
quantity, not $\beta$; session reset ($k=0$ per probe) is the identification
condition.
\end{proof}

\begin{proof}[Proof of Proposition~\ref{prop:alloc}]\looseness=-1
Arithmetic on the hide threshold. A probe-aware target hides on a family once
its exposure there reaches $K$. If the family is probed $n_f<K$ times, the
measured rate saturates at $\hat\beta_f=\min(1,K/n_f)=1$, and the family fails
any tolerance. If $n_f\ge K$, then $\hat\beta_f=K/n_f$, and the certificate at
tolerance $\tau$ fails on that family exactly when $K/n_f>\tau$, i.e.\
$n_f<K/\tau$. Uniform allocation hands the attacker the uniform
exposure it needs to hide everywhere; concentrated allocation leaves some
families under-exposed, where it cannot hide. No new inequality is claimed;
the content is the operational reversal.
\end{proof}

\begin{proof}[Proof of Proposition~\ref{prop:recon}]
Decision analysis under the published vector; finite-sample uncertainty is
reported via the Wilson bounds---the proposition registers the rule.
\end{proof}

\begin{proof}[Proof of Proposition~\ref{prop:release}]\looseness=-1
Part (i) is Theorem~\ref{thm:forgery}'s realizability with the defender as
author; Part (ii) is Proposition~\ref{prop:recon}'s allocation analysis
with a defender-chosen argument; Part (iii) reverses
Corollary~\ref{cor:adaptive}'s secrecy asymmetry. No new inequality is
claimed; the new objects are the release gain (honest vs.\ decoy success)
and the verification-stability condition.
\end{proof}

\begin{proof}[Proof of Proposition~\ref{prop:budget}]\looseness=-1
The public support carries no positives, so any public-side misread fails
the certificate immediately: the one-sided asymmetry is the content, and
the public-side misread rate $m'$ must be zero for detection of public
breaches at all. The attacker therefore has exactly one detection channel,
the held-out stream: a breached probe counts as a breach iff the oracle
classifies it as held-out, so the detected-breach count is
$\mathrm{Bin}(n,1-m)$; the detection power is the binomial tail at the
first count $k_{\min}$ whose Wilson lower bound exceeds the public
certificate's upper bound (0.0602 for $0/60$), and $n_{\min}$ is the
smallest $n$ with that tail at 95\%. As $m\to1$ the per-probe detection
probability $1-m\to0$ and the tail needs a larger $n$ at fixed $\alpha$
until no finite $n$ suffices: evasion requires error-free one-sided
discrimination of the public support, and no finite budget guarantees
detection at $\alpha<1$.
\end{proof}
